% Proposed insertion, not applied to the repository.
% Requires the existing theorem/corollary environments and amsmath/amssymb.
% No global shorthand macros are introduced. Full proofs are in the report
% integral-distribution-reflection/article/article.pdf, Sections 2--6.
\section{Integral distribution, reflection torsion, and modular jets}
\label{ir:insert:section}

Let $X_q=(q^{-1}\mathbb Z)/\mathbb Z$ and
$R_q=\mathbb Z[A_p:p\mid q]$. Define
\[
 \mathcal D_q=\left(\bigoplus_{x\in X_q}R_qe_x\right)
 \bigg/\left\langle\sum_{py=x}e_y-A_pe_x:
          p\mid q\text{ prime},\ x\in X_{q/p}\right\rangle.
\]
The point $0$ is formal; its Hurwitz evaluation is at $1$, not $0$.
The characteristic-zero conductor basis can be supplemented by a
basis requiring no character projectors or divisions. Ordinary
specializations have classical antecedents
\cite{ir:KubertSign1979,ir:Ouyang2002}; the companion report supplies
self-contained proofs for this weighted presentation and its deformations.

\begin{theorem}[Division-free distribution basis]
\label{ir:insert:integral}
For each $p^a$, discard one lift from each fiber of
$(\mathbb Z/p^a\mathbb Z)^\times\to
 (\mathbb Z/p^{a-1}\mathbb Z)^\times$, using a one-element set at
exponent zero. Retain $0$ and all rational grid points whose nonzero
primary coordinates avoid the discarded lifts. Their images form an
$R_q$-basis of $\mathcal D_q$, of cardinality $\varphi(q)$.
The same basis works after every commutative-ring specialization.
The raw prime-relation matrix has a selected maximal minor equal to
$1$, of size $q-\varphi(q)$.
\end{theorem}
The terminating rule is
\[
 e_y=A_pe_{py}-\sum_{pz=py,\ z\ne y}e_z
\]
at a discarded $p$-primary lift. Order points by denominator and then
by number of discarded components. This rule strictly decreases that
order. Independence follows from the finite resolution with chain
groups $\bigoplus_{|S|=j}R_q[X_{q/\prod_{p\in S}p}]$: filtration by
$\operatorname{den}(x)\prod_{p\in S}p$ has associated graded complex
a tensor product of split local norm injections. Thus finite rank
experiments are replaced by an all-ring proof.

Let $c(e_x)=e_{-x}$. For $q>2$ put
\[
 n=\varphi(q)/2,\qquad
 r=\#\{p>2:p\mid q\}+\mathbf1_{4\mid q},\qquad h=2^{r-1}.
\]
\begin{theorem}[Integer reflection coinvariants]
\label{ir:insert:integer-reflection}
For integer weights $a_p$, and either sign $\varepsilon=\pm1$,
\[
 \mathcal D_q(\mathbb Z,a)/(c-\varepsilon)\mathcal D_q(\mathbb Z,a)
 \simeq\mathbb Z^n\oplus(\mathbb Z/2\mathbb Z)^t,
\]
where $t=h$ if every odd-prime weight is odd, and $t=0$ otherwise.
The answer does not depend on the integer $a_2$.
At $q=1,2$, the plus and minus answers are respectively
$\mathbb Z$ and $\mathbb Z/2\mathbb Z$.
\end{theorem}
The unreflected relation module is torsion-free; this theorem concerns
a subsequent symmetry quotient. Its proof removes the moving two-point
orbits from the periodic cohomology of the finite resolution. The fixed
complex reduces to the Koszul complex on
\[
 (a_p-1:p>2,\ p\mid q),
 \quad\text{with }a_2(a_2-1)\text{ added when }4\mid q.
\]
For $v_2(q)=1$ the fixed-point column at $2$ is
$(1-a_2,1)^{\mathsf T}$ and contributes no extra Koszul direction.
For $v_2(q)\ge2$ its matrix is
$\left(\begin{smallmatrix}1-a_2&0\\1&-a_2\end{smallmatrix}\right)$,
which has a unit pivot and determinant $a_2(a_2-1)$.

\begin{theorem}[Characteristic-two contact-order law]
\label{ir:insert:jets}
Let $k$ have characteristic two, take weights $a_p(t)\in k[[t]]$,
and let $d$ be the least $t$-adic valuation of the active scalars just
listed, with $v_t(0)=\infty$. For finite $d$, the Smith diagonal of
$c-1$ on $\mathcal D_q(k[[t]],a)$ consists of $n-h$ units,
$h$ entries $t^d$ up to units, and $n$ zeros. For $d=\infty$ it has
$n-h$ units and $n+h$ zeros. Consequently for all $L\ge1$,
\[
 \dim_k\frac{\mathcal D_q(k[t]/(t^L),a)}
 {(c-1)\mathcal D_q(k[t]/(t^L),a)}
 =nL+h\min(d,L).
\]
\end{theorem}
More generally, over every characteristic-two algebra $S$,
$\ker(c-1)/\operatorname{im}(c-1)$ is the direct sum of the homology
groups of the active Koszul complex. Over the universal polynomial
ring it is the cyclic quotient by the active ideal. Over a field the
reflection dimension is $n+h$ precisely when all odd-prime weights
are $1$ and, if $4\mid q$, $a_2$ is $0$ or $1$; otherwise it is $n$.
These modular jets are formal weight deformations, not reductions
modulo two of complex logarithms.

\subsection{A short all-order analytic certificate}
Put $\mathcal L_s(x)=\operatorname{Li}_s(e^{2\pi i x})$ and
$U=\{1/15,4/15,7/15,2/5,4/5\}$. The four raw rows
\[
 r_{3,3/5}-r_{5,1/3}-A_5r_{3,0}-r_{5,0}
\]
prove, after $A_p=p^{1-s}$,
\begin{align*}
 &\mathcal L_s(8/15)-\sum_{u\in U}\mathcal L_s(u)
 -(3^{1-s}+1)\mathcal L_s(3/5)\\
 &\qquad -(5^{1-s}+1)\mathcal L_s(2/3)
 =(1-15^{1-s})\zeta(s).
\end{align*}
Both sides extend holomorphically through $s=1$ with right-hand value
$\log15$. If $\ell=\log15$, its $N$th derivative there has scalar
right-hand side
\[
 (-1)^N\left(\frac{\ell^{N+1}}{N+1}
 -\sum_{j=1}^N\binom Nj\ell^j\gamma_{N-j}\right),\qquad N\ge0,
\]
with ordinary Leibniz derivatives of the weighted terms on the left.
The term $\ell^{N+1}/(N+1)$ is the contribution of the zeta pole and
must not be discarded before differentiation.

These are formal-module and exact-identity theorems. They do not
establish numerical independence of periods or prove the separate
mixed Gaussian $S_6$ conjecture. The companion report includes the
complete proofs, frozen JSON certificates, and exact replay code.
